\section*{Bab \ref{ch:g11:lithosphere-asthenosphere} --- Litosfer dan Astenosfer}

\begin{solution}{exo:g11:lithosphere-asthenosphere:1}
\omterm{def:g11:lithosphere-asthenosphere:lithosphere}{Litosfer} adalah kulit luar Bumi yang kaku, yang terdiri atas kerak dan mantel
paling atas; \omterm{def:g11:lithosphere-asthenosphere:lithosphere}{astenosfer} adalah mantel yang lebih panas dan lebih lemah di
bawahnya, yang mengalir perlahan dalam waktu yang lama. Kerak samudra, kerak
benua, dan mantel paling atas termasuk \omterm{def:g11:lithosphere-asthenosphere:lithosphere}{litosfer}; inti luar tidak.
\end{solution}

\begin{solution}{exo:g11:lithosphere-asthenosphere:2}
Gunung es itu tenggelam sampai mendesak air yang beratnya sama dengan beratnya
sendiri: bagian di bawah air adalah $910/1025 \approx 0.888$, sekitar
\qty{89}{\percent}. Bagian di atasnya sekitar \qty{11.2}{\percent}, jadi tinggi
totalnya $20/0.112 \approx \qty{180}{m}$, dengan sekitar \qty{160}{m} di bawah
air.
\end{solution}

\begin{solution}{exo:g11:lithosphere-asthenosphere:3}
$15 + 26 \times 3 = \qty{93}{\celsius}$. Pada \qty{100}{km} gradien yang sama
akan memberikan lebih dari \qty{2600}{\celsius}, jauh di atas suhu leleh batuan
mantel; padahal \omterm{def:g11:seismic-structure:body-waves}{gelombang S} menyeberangi mantel atas, yang karena itu padat.
Gradiennya harus menjadi jauh lebih kecil seiring kedalaman.
\end{solution}

\begin{solution}{exo:g11:lithosphere-asthenosphere:4}
Pita dari 0 sampai \qty{500}{m} di atas permukaan laut, dengan sekitar
\qty{16}{\percent} permukaan, dan pita dari 4000 sampai \qty{4500}{m} di bawah
permukaan laut, dengan sekitar \qty{11}{\percent} (pita dari 4500 sampai
\qty{5000}{m} di bawahnya hampir sama besarnya).
\end{solution}

\begin{solution}{exo:g11:lithosphere-asthenosphere:5}
Pada $P = 0$: \qty{1121}{\celsius}. Pada \qty{3}{GPa}: $-5.104 \times 9 +
132.899 \times 3 + 1120.661 \approx -45.9 + 398.7 + 1120.7 \approx
\qty{1473}{\celsius}$.
\end{solution}

\begin{solution}{exo:g11:lithosphere-asthenosphere:6}
Sekitar \qty{730}{\celsius} pada \qty{50}{km} (\omterm{def:g11:lithosphere-asthenosphere:lithosphere}{litosfer}), sekitar
\qty{1370}{\celsius} pada \qty{150}{km} dan \qty{1410}{\celsius} pada
\qty{250}{km} (\omterm{def:g11:lithosphere-asthenosphere:lithosphere}{astenosfer}). Suhu naik paling banyak antara 50 dan
\qty{150}{km}, lebih dari \qty{600}{\celsius}, karena sebagian besar selang itu
terletak di \omterm{def:g11:lithosphere-asthenosphere:lithosphere}{litosfer}; antara 150 dan \qty{250}{km} suhu hanya naik sekitar
\qty{45}{\celsius}.
\end{solution}

\begin{solution}{exo:g11:lithosphere-asthenosphere:7}
Akar: $3 \times 2830/550 \approx \qty{15.4}{km}$. Tambahan tebal:
$3 \times 3380/550 \approx \qty{18.4}{km}$, yaitu \qty{3}{km} jajaran
pegunungan itu ditambah akarnya.
\end{solution}

\begin{solution}{exo:g11:lithosphere-asthenosphere:8}
Kerak lebih tebal sebesar $55 - 41.1 = \qty{13.9}{km}$, dan $\Delta H = h
\times 3380/550$, jadi $h = 13.9 \times 550/3380 \approx \qty{2.3}{km}$.
\end{solution}

\begin{solution}{exo:g11:lithosphere-asthenosphere:9}
Penurunan: $(4.470 - 4.419)/4.470 \approx \qty{1.1}{\percent}$. Lonjakan:
$(4.644 - 4.419)/4.419 \approx \qty{5.1}{\percent}$. Seiring kedalaman batuan
makin lama makin termampatkan, yang membuat gelombang seismik lebih cepat;
penurunan berarti ada hal lain yang menang: batuan itu panas dan hampir meleleh,
dan lebih lemah.
\end{solution}

\begin{solution}{exo:g11:lithosphere-asthenosphere:10}
(1) $(1345 - 610)/60 \approx \qty{12}{\celsius}$ per kilometer. (2) Nilai itu
sekitar setengah gradien di permukaan. Sebagian kalor dihasilkan di dalam kerak
atas; kalor itu juga harus menyeberangi lapisan-lapisan atas untuk keluar, yang
membuat gradien lebih curam di dekat permukaan. Di bawah kerak hanya kalor yang
datang dari bagian yang lebih dalam yang menyeberangi batuan, dan gradiennya
lebih kecil.
\end{solution}

\begin{solution}{exo:g11:lithosphere-asthenosphere:11}
Pada \qty{200}{km}: $1300 + 90 = \qty{1390}{\celsius}$; pada \qty{400}{km}:
$1300 + 180 = \qty{1480}{\celsius}$. Kenaikannya \qty{90}{\celsius} sepanjang
\qty{200}{km}. Melalui \omterm{def:g11:lithosphere-asthenosphere:lithosphere}{litosfer} suhu naik sebesar
$1345 - 15 = \qty{1330}{\celsius}$ sepanjang hanya \qty{100}{km}: sekitar tiga
puluh kali lebih banyak per kilometer.
\end{solution}

\begin{solution}{exo:g11:lithosphere-asthenosphere:12}
(1) $118 + 12.3\,z = 1300 + 0.45\,z$ memberikan $z = 1182/11.85 \approx
\qty{100}{km}$. (2) $118 + 21.5\,z = 1300 + 0.45\,z$ memberikan $z = 1182/21.05
\approx \qty{56}{km}$. (3) Bagian barat yang lebih panas mempunyai \omterm{def:g11:lithosphere-asthenosphere:lithosphere}{litosfer} yang
jauh lebih tipis, sesuai dengan pengukuran (sekitar \qty{40}{km} di tempat yang
paling tipis); \omterm{def:g11:lithosphere-asthenosphere:lithosphere}{litosfer} yang sangat tebal di bagian tengah benua itu memerlukan
gradien yang lebih kecil lagi.
\end{solution}

\begin{solution}{exo:g11:lithosphere-asthenosphere:13}
(1) Menurut proposisi itu, $\Delta H = h\,\rho_m/(\rho_m - \rho_c)$, jadi $h =
\Delta H\,(\rho_m - \rho_c)/\rho_m$. (2) $550/3380 \approx 0.16$: sekitar
seperenam dari setiap tambahan kerak tampak sebagai ketinggian, lima perenam
sebagai akar. (3) Untuk gunung es, bagian di atas air adalah
$(1025 - 910)/1025 \approx 0.11$: aturannya sama, (massa jenis fluida $-$ massa
jenis benda yang mengapung)/massa jenis fluida.
\end{solution}

\begin{solution}{exo:g11:lithosphere-asthenosphere:14}
(1) $1489 - 1345 = \qty{144}{\celsius}$. (2) Sedikit lelehan berkumpul dalam
selaput tipis di antara butir-butir batuan. Butir-butir itu kemudian tidak lagi
terkunci erat satu sama lain: batuan itu kurang melawan geseran, sehingga
\omterm{def:g11:seismic-structure:body-waves}{gelombang S} melambat, dan batuan itu dapat berubah bentuk dan mengalir jauh
lebih mudah.
\end{solution}

\begin{solution}{exo:g11:lithosphere-asthenosphere:15}
(1) Tidak. Meratakan daratan akan membuat kerak benua sedikit lebih tipis,
tetapi kerak itu tetap tebal dan ringan, dan tetap mengapung jauh lebih tinggi
daripada kerak dasar samudra yang tipis dan rapat; samudra tetap akan mengisi
cekungan di antara benua-benua. (2) Ada dua jenis kerak, dengan tebal dan massa
jenis yang berbeda: kerak benua yang tebal dan ringan dan kerak samudra yang
tipis dan rapat.
\end{solution}

\begin{solution}{pb:g11:lithosphere-asthenosphere:1}
\textbf{1.} $5.01 - 0.67 = \qty{4.34}{km}$.

\textbf{2.} $(3380 - 2830)/3380 \approx \qty{16}{\percent}$.

\textbf{3.} Bagian $2830/3380 \approx 0.84$ akan terletak di bawah; gunung es
tenggelam sedikit lebih dalam, $910/1025 \approx 0.89$, karena massa jenis es
lebih dekat dengan air daripada massa jenis kerak dengan mantel.

\textbf{4.} Selama beberapa detik perambatan gelombang seismik, atau selama
gempa bumi, mantel adalah zat padat; hanya dalam waktu yang sangat lama
\omterm{def:g11:lithosphere-asthenosphere:lithosphere}{astenosfer} mengalir dan membuat kerak mengapung.

\textbf{5.} Di bawah \omterm{def:g11:seismic-structure:moho}{Moho} yang paling dalam, kedua kolom terbuat dari batuan
mantel yang sama pada kedalaman yang sama, dan di sana beratnya sama; semua
perbedaan terletak di atas kedalaman itu dan harus diperhitungkan.

\textbf{6.} Benua rata-rata: kerak dari \qty{0.67}{km} di atas sampai
\qty{40.4}{km} di bawah permukaan laut, lalu mantel sampai \qty{70}{km}. Dataran
tinggi: kerak dari \qty{5.01}{km} di atas permukaan laut sampai \omterm{def:g11:seismic-structure:moho}{Moho} yang lebih
dalam sebesar akarnya, lalu mantel sampai \qty{70}{km}.

\textbf{7.} Kolom dataran tinggi mempunyai kerak $h + r$ lebih banyak, sedangkan
kolom benua mempunyai mantel setebal $r$ sebagai gantinya: $\rho_c\,(h + r) = \rho_m\,r$,
jadi $r = h\,\rho_c/(\rho_m - \rho_c)$ dan $\Delta H = h + r =
h\,\rho_m/(\rho_m - \rho_c)$.

\textbf{8.} $3380/550 \approx 6.15$.

\textbf{9.} $\Delta H = 4.34 \times 6.15 \approx \qty{26.7}{km}$; akarnya
$26.7 - 4.34 \approx \qty{22.3}{km}$.

\textbf{10.} \omterm{def:g11:seismic-structure:moho}{Moho} benua rata-rata terletak $41.1 - 0.67 \approx
\qty{40.4}{km}$ di bawah permukaan laut; \omterm{def:g11:seismic-structure:moho}{Moho} dataran tinggi terletak $40.4 + 22.3 \approx
\qty{62.8}{km}$ di bawah permukaan laut.

\textbf{11.} $41.1 + 26.7 \approx \qty{67.8}{km}$.

\textbf{12.} Satu kilometer batuan di atas satu meter persegi bermassa
$\rho \times \qty{1000}{m^3}$: \num{2830} ton kerak atau \num{3380}
ton mantel. Benua: $41.1 \times 2830 + (70 - 40.43) \times 3380
\approx \num{116300} + \num{99900} \approx \num{216300}$ ton. Dataran tinggi:
$67.77 \times 2830 + (70 - 62.76) \times 3380 \approx \num{191800} +
\num{24500} \approx \num{216300}$ ton. Keduanya sama.

\textbf{13.} $41.1 - 0.67 \times 6.15 \approx \qty{37.0}{km}$.

\textbf{14.} $67.8/41.1 \approx 1.65$ kali.

\textbf{15.} $3300/470 \approx 7.02$, $\Delta H \approx 4.34 \times 7.02
\approx \qty{30.5}{km}$, dan kerak setebal sekitar \qty{71.6}{km}. Perubahan
massa jenis mantel sekitar \qty{2}{\percent} mengubah tebal itu sekitar
\qty{4}{km}: hasilnya cukup bergantung pada massa jenis, tetapi tetap dekat
dengan \qty{70}{km}.

\textbf{16.} Model itu memberikan sekitar \qty{68}{km}, gelombang seismik sekitar
\qty{70}{km}: keduanya sesuai sampai beberapa persen.

\textbf{17.} $h = 3.23 - 0.67 = \qty{2.56}{km}$, $\Delta H = 2.56 \times
6.15 \approx \qty{15.7}{km}$, dan kerak setebal sekitar \qty{56.8}{km}, sekitar
\qty{7}{km} lebih tebal daripada yang terukur.

\textbf{18.} Lebih tipis. Kolom yang bagian atasnya lebih ringan daripada
rata-rata lebih ringan untuk tebal yang sama, sehingga mengapung lebih tinggi:
kerak yang lebih ringan memerlukan akar yang lebih kecil untuk menjulang pada
ketinggian tertentu. Karena itu kerak yang sebenarnya dapat lebih tipis daripada
kerak model, sesuai dengan pengukuran.

\textbf{19.} Dalam waktu yang lama \omterm{def:g11:lithosphere-asthenosphere:lithosphere}{astenosfer} mengalir: \omterm{def:g11:lithosphere-asthenosphere:lithosphere}{astenosfer} menyingkir di
bawah kolom yang berat dan kembali di bawah kolom yang ringan, sampai kolom-kolom
itu seimbang.

\textbf{20.} \textbf{Kerak di bawah Tibet setebal sekitar \qty{70}{km}}, jauh
lebih tebal daripada \qty{41}{km} benua rata-rata: dataran tinggi itu menjulang
lima kilometer karena keraknya yang tebal dan ringan mengapung di atas mantel
yang lebih rapat, dengan akar lebih dari \qty{20}{km}.
\end{solution}
